\section{Community of Blood vs Individualism}

\textbf{Rudolf Steiner} delivered a lecture on 15 June 1915\footnote{\url{https://rsarchive.org/Lectures/PreSix_index.html}} that has become relevant to world events happening in our own day. This is how he starts:

\begin{quotex}
It is our desire and aim to recognize out of knowledge that the world has meaning, significance and purpose, and that the world is not filled merely with evil and degeneration. It is our aim to realize through direct knowledge that the world has meaning. By this realization we try to prepare for actual experience of the Christ.
\end{quotex}


\textbf{Oswald Spengler} was working on his two volume masterpiece, \emph{The Decline of the West}, around the same time as Steiner's lectures. Both were heavily influenced by \textbf{Friedrich Nietzsche} and \textbf{Johann Wolfgang von Goethe}, and both were interested in the cycles of cultures. Despite their differences in approach and world outlook, there are certainly correspondences between the two.

Unlike Spengler, Steiner was not interested in isolated cultures like the Mexican. Rather, he recognized a continuity as the end of one culture gives birth to the next. Steiner's cultures, or ages, are then close to Spengler's, as shown below:

\begin{enumerate}


\item Ancient India

\item Persian

\item Egyptian

\item Greco-Latin

\item Western Europe (and its diaspora in North America and Australia). This is the current dominant culture.

\item To come: East Europe, Slavic/Russian
\end{enumerate}


This matches up, at least in general outline, with Spengler's list of cultures. However, there is one significant difference between them. Spengler believed that the process of the birth, decline, and death of a culture was blind, unplanned, and unrelated to each other. He explains:

\begin{quotex}
Mankind... has no aim, no idea, no plan, any more than the family of butterflies or orchids. `Mankind' is a zoological expression, or an empty word. ... I see, in place of that empty figment of one linear history which can only be kept up by shutting one's eyes to the overwhelming multitude of the facts, the drama of a number of mighty Cultures, each springing with primitive strength from the soil of a mother region to which it remains firmly bound throughout its whole life-cycle; each stamping its material, its mankind, in its own image; each having its own idea, its own passions, its own life, will and feeling, its own death.
\end{quotex}


Steiner, on the other hand, saw the development of successive cultures as part of a larger process. Rather than a blind process, each successive stage was prepared by a spiritual elite. In brief, he wrote:

\begin{quotex}
The centers of the mysteries were the places in which the form of external life belonging to the next epoch of culture was prepared. The mysteries were associations of human beings among whom other things were cultivated than those cultivated in the outer world.
\end{quotex}


Hence, there is an overlap during the transition from one cultural age to the next. The psychic and spiritual environment has to be prepared for souls with affinity to such environments to come to birth in each age. One thing is certain, however: those at one stage can only with difficulty understand those at another.

\paragraph{Characteristics}


\subparagraph{Apollonian}


Spengler referred to the Classical Greco-Latin culture as Apollonian. It was characterized by its preference for the local and the present moment.

This is similar to Steiner's conception. The ancient Greek identified with his city; an inhabitant of Athens saw himself as an ``Athenian''. Similar for the citizen of Rome. Outside of the city, there were only barbarians. Cosmopolitanism was an alien concept because, at that age, a human being was a member of a group.

\subparagraph{Faustian}


Spengler believed that this Faustian culture is characterized by the unlimited search for knowledge, as exemplified by the elusive search for the ``Theory of Everything''. He figured that it is has run its course since the goal is unachievable and predicted that it would begin its final decline right around the current time.

Steiner said that, at this age, the individual strives to grow beyond the community, that is, to be a ``human being'', pure and simple, transcending his local community ties. This is expressed by extreme individualism, with each man his own master. But individualism is the principle of multiplicity, so this age led to several major conflicts.

\subparagraph{Slavic Age}


Both Spengler and Steiner saw a special role for Russia. Spengler predicted a new High Culture emerging from Russia, which Steiner called the Sixth Age. However, Russia is still at the Greco-Roman stage since the Russian soul still identifies with the community. They regard this as their birthright and achievement and are therefore opposed to the extreme individualism or the West. That attitude is retrograde, if not dangerous, and its true destiny is to establish communities based on spiritual affinity rather than on blood.

\paragraph{World War X}


The war that was being fought in 1915 is analogous to the European conflict now. Then, the instigation for the war was Austria vs Russia/Slavs, representing the fifth cultural age and the fourth, respectively. (The war was triggered in Serbia---aligned with Russia---against Austria). Steiner defines the former:

\begin{quotex}
A great and terrible symbol stands before the eyes of the world. Think of the two states where the war had its starting-point. On the one side, Russia with the Slavic world in general, declares that the war is based on brotherhood of blood, and on the other side, there is Austria, which comprises thirteen distinct peoples and thirteen different languages. The mobilization order in Austria had to be issued in thirteen languages because Austria encompasses thirteen racial stocks: Germans, Czechs, Poles, Ruthenians, Rumanians, Magyars, Slovaks, Serbs, Croatians, Slovenes (among whom there is a second and separate dialect), Bosnians, Dalmatians and Italians. Thirteen different racial stocks, apart from all minor differentiations, are united in Austria. Whether the implications of this are understood or not, it is obvious that Austria consists of a collection of human beings among whom community can never be based on blood relationship, for what its strange boundaries contain shoots out into thirteen different lineages.
\end{quotex}


This was opposed to the Russian soul:

\begin{quotex}
The most highly composite state in Europe stands in opposition to the state that strives most intensively for life in a group soul, or for conformity. But this striving for life in a group soul brings a great many other things in its train.
\end{quotex}


The analog to today is irresistible. Now the EU and NATO, represent the polyglot raceless state, against the Russian soul based on blood. This is confirmed by recent talks delivered both by Putin and by the Patriarch of the Russian Church. Putin believes, as he made clear, that Ukraine is still part of the Slavic community of blood. But Ukraine has been emerging from that stage as the attraction of the Faustian stage is strong.

There seems to be little possibility of any mutual understanding between the Russian ideal of the community of blood and the Western ideal of the cosmopolitanism of individuals. The West regards Putin as mad and Putin regards the West as evil. The Russian church is adamantly opposed to Western ``rights'' like same sex marriage and transgenderism. That baffles the West, which has no sympathy at all for the Patriarch's worldview.

\paragraph{Common Task}


To remain at the community of blood, while a necessary stage at the time, is no longer appropriate. Individualism leads to atomism, thereby inhibiting common action unless imposed by the forces of ideology. It has no explanation for the origins or eradication of evil; neither new government programs, nor alterations to the social structure have been adequate. Evil persists and unhappiness endures in the West.

The Common Task inaugurated by Nikolai Fyodorov, is a start. Intentional communities of spirit will replace dying cultural relics. The Christian world has yet to fully understand the process of the streams of thought inaugurated by Vladmir Solovyov and continued by Sergei Bulgakov and Nikolai Berdyaev.

Then there is the esoteric current in the Slavic soul that has been brought into the daylight by Valentin Tomberg. He recapitulates ideas from G. O. Mebes, Nina Roudnikova, and Vladimir Schmakov, which have only recently been published in English.

Steiner's view is decidedly Eurocentric. Although it recognizes the contribution of the Vedas to the esoteric tradition, he leaves out the contributions of the Islamic World, or Magian in Spengler's understanding. For this, we need to rely on the works of Rene Guenon.

{Part 2, on Evil in the World, will appear tomorrow.}


\flrightit{Posted on 2022-03-12 by Cologero}

\begin{center}* * *\end{center}

\begin{footnotesize}
\begin{sffamily}
\texttt{Imperius on 2022-03-13 at 00:12 said:}

Russia is itself a multiracial empire, though. And much of the country is against the invasion, already converged by Western capitalism, not crying out for national identity.

Truly, at this point Putin will be lucky if he is not assassinated, all the collapse in trade.

And by the way, northern Russian cities are more cosmopolitan than Ukraine. Putin is really just a pragmatist, not some avatar.

\hfill

\texttt{Cologero on 2022-03-13 at 09:30 said:}

It's a possibility that you totally missed the point.The temperature and pressure of a gas can be measured without knowing the positions and motions of each individual atom. That is because the gas is a totality, so the differences in individual atoms turn out to be irrelevant because they average out. Moreover, the individual atoms are not even aware that they are participants in a larger process.

Granted, the cyclic view of history, like that of Spengler or Steiner, is controversial. Perhaps things just happen without purpose or meaning.

\hfill

\texttt{Arthur Konrad on 2022-03-13 at 11:02 said:}

@ Imperius

What anybody in Russia thinks is totally unimportant, because democracy never existed in Russia other than this Potemkin democracy they orchestrate for whatever reason. Until the collapse of the Russian Empire there has never been any sort of representation in the Russian political system, even less in the Soviet Union. This tradition carried on until the present. It appears you are trying to argue that Russia would be a better place if `much of the country' had a say in the way it is governed, and not even more degenerate than the modern West.

There is a saying `misery loves company' which describes very well the kind of anguish the West experiences when dealing with self-confident cultures. Westerners would like you to believe that their loss of prestige and importance had something to do with noble reasons like justice, equality, etc, and not with its foolish choice to abolish all standards and concern for the future in order to enjoy the present in the mist of cannabis smoke and guitar riffs. So subconsciously, they would like everyone else to descend to their level. People like Putin, Jinping or Prince Salman will have none of that. They understand very well that the West lost its place in the world because it is weak and pathetic, and hence cannot fail to express their contempt for Western values at every opportunity. In this sense the East is `nationalistic', however much this definition of `nationalism', or `statism' does not conform to one or another flexible and totally unprincipled schema for winning debates about why your enemies shouldn't exist.

\hfill

\texttt{Kris on 2022-03-13 at 13:07 said:}

Rootless cosmopolitanism is the handmaiden of liberal plutocratic capitalism. Many Russians see through the intellectual bankruptcy and spiritual emptiness of Western liberal democracy' and rootless cosmopolitanism bus mass immigration, feminism, egalitarianism, and other evils.

\hfill

\texttt{Balder on 2022-03-14 at 05:56 said:}

This seems most close to an elaboration of current events that I hope for some time ago. Very good!

I'm not an expert either with Spengler or Steiner but I guess they both had some insights although Steiner seemed to elevate rootless individualism. Spengler seemed to be a fatalist; can't really blame him for that.

I can't support Putin in any way in his offensive, but while many if not most American companies (McDonald's, Coca-Cola etc.) are leaving the country, western ``values'' are downgraded and traditionalist Orthodoxy seems to be venerated (still considering conversion), I think I'd like living in Russia in the near future. Perhaps I'd have the fate of a useful idiot there still because of a wrong nationality.

\hfill

\texttt{Greg on 2022-03-14 at 13:07 said:}

Fyodorovich and his common task of resurrecting the dead and colonizing the world through science and Christianity sounds more in line with the transhumanism of the WEF then Tradition. What am I missing? My first impulse us that you were joking.

\hfill

\texttt{Balrog Booger on 2022-03-14 at 13:46 said:}

For anyone with about 45 minutes to spare, this talk by Branko Malic, the Croation proprietor of the Tradition-adjacent blog \emph{Kali Tribune,} is interesting. I'm a typical American who is totally ignorant about many aspects of Eastern European history, and his explanation of Russia's tripartite sense of being the \emph{Third Rome} is something I have never been exposed to before. If his view has any truth to it, it explains quite well the cultural impetus that could lead Russia to be the next dominant world power.

\url{http://en.kalitribune.com/arrivederci-third-rome/}


\hfill

\texttt{Cologero on 2022-03-14 at 14:05 said:}

Sort of joking, Greg, or just curious to see if anyone would notice. My primary source is ``The Russian Cosmists'' by George Young although I've read extracts from Fyodorov's book. I don't understand his influence on thinkers I otherwise respect.

A good friend has been translating G. O. Mebes, Nina Roudnikova, and Vladimir Schmakov. When the Schmakov book is ready, I told here I would do a composite review.

\hfill

\texttt{Max on 2022-03-14 at 14:35 said:}

The Slavic world grows to despise Russian actions more for every day that passes. They are currently at the forefront of establishing intentional communities of spirit, precisely by rejecting the retrograde Russian culture, as expressed here. Soon there might be nobody left who would voluntarily want to be part of it. Those who nevertheless still do, should be sent on their way. That is how intentional community is formed in practice. Back in the days, there were even people who left for the soviet union, and reportedly had a great time there.

Individualism is said to have ``no explanation for the origins or eradication of evil''. That is not correct, for it sees evil as being in each individual man. The ``community of blood'' on the other hand, tends to engage in genocidal wars when it somehow feels wronged.

As to capability for collective action, the last chapter has not been written yet, although the initial inertia of atomistic indecision was definitely overcome. It is noteworthy that Western Slavic nations have taken up leadership roles, as they should. They should know more than any of what is at stake.

\hfill

\texttt{Adil on 2022-03-14 at 15:42 said:}

We are returning to the bi-partisan world of the Cold War between communism and capitalism, but this time the dichotomy is more real -- that between liberalism and tradition. The Cold War was just a pretext to a deeper more metaphysical rift. It is interesting that Russia now has both Orthodox and Muslims fighting on the same side against the West.

\hfill

\texttt{Logres on 2022-03-19 at 10:28 said:}

An interesting twist of history would be that, should the West prove too incompetent and corrupt on a vast scale, Russia could actually triumph and humiliate the Faustian order. This would give temporary validation to the old ethnic impulses and further complicate matters (similar to how America humiliated Britain in the colonial era). It's by no means clear that they won't succeed in doing so, since our purblind choices here have made that more of a possibility. There is lots of potential to re-arrange historical Karma. Russia asserted a limit. Thank you, a yeoman like essay. Very Sun-Tzu and also, percipient.

\end{sffamily}
\end{footnotesize}

